\chapter*{Resumen}
\addcontentsline{toc}{chapter}{Resumen}

Este Manual Tecnico documenta la Plataforma Web Robot Explota Globos, sistema desarrollado para soportar el registro de participantes, la gestion de comunicaciones y el control operativo del torneo Robot Explota Globos. El proyecto pertenece al contexto de la Universidad Tecnologica de Pereira y del Semillero de Investigacion Kilby. La responsabilidad del proyecto esta asociada a Yurley Tatiana Tovar Martinez, y el desarrollo fue realizado por Moises Sanchez Naranjo y Tomas Bermudez Londono.

La plataforma esta implementada como un proyecto Django modular. La auditoria tecnica validada identifico cinco aplicaciones locales: \archivo{apps.core}, \archivo{apps.participants}, \archivo{apps.tournament}, \archivo{apps.invitations} y \archivo{apps.common}. Esta separacion responde a responsabilidades distintas: sitio publico, inscripciones y asistencia, logica de competencia, comunicaciones institucionales y reutilizacion de modelos base. La estructura permite que los flujos publicos de registro no queden mezclados con las operaciones internas de torneo y comunicaciones, lo que mejora la mantenibilidad del sistema y reduce el riesgo de cambios cruzados entre dominios.

El componente de competencia es el nucleo tecnico mas sensible. El sistema soporta registros de colegio y universidad, confirmacion de asistencia mediante token UUID, sincronizacion de registros confirmados hacia equipos operativos y construccion de competencias por division. La logica de torneo contempla fase de grupos, eliminatorias, repechajes, fases progresivas, fases manuales, correccion de resultados y registro de podio final. Esta capacidad esta centralizada principalmente en \archivo{apps/tournament/services.py}, decision que facilita ubicar la regla de negocio principal, aunque tambien concentra riesgo de regresion si se modifica sin pruebas suficientes.

El modulo de comunicaciones gestiona plantillas DOCX, destinatarios XLSX, renderizado de documentos, conversion a PDF mediante LibreOffice y envio de correos por backend SMTP o consola, segun configuracion. Este modulo incorpora modos de simulacion y prueba controlada para reducir el riesgo de envios reales accidentales. El frontend se mantiene en Django Templates, CSS compartido y JavaScript especifico: \archivo{static/js/app.js} para interacciones publicas y \archivo{static/js/control.js} para el panel interno, incluyendo operaciones AJAX con CSRF.

La auditoria tecnica fue calificada con 88/100 y quedo aprobada con observaciones. Las observaciones no bloquean la redaccion del manual, pero orientan los capitulos tecnicos posteriores: detallar el esquema de \variable{DivisionCompetition.configuration}, ampliar contratos internos de servicios criticos, convertir brechas de prueba en matriz de QA, no inventar datos de despliegue y mantener pendientes de produccion marcados como \MarcadorProduccionPendiente. Este primer bloque establece el contexto, objetivos, alcance, descripcion general y arquitectura implementada que serviran como base para los capitulos tecnicos restantes.
